% =====================================================================
% Fluid Dynamics and IBM Theory: immersed boundary method discretization
% =====================================================================
\documentclass[aspectratio=169,11pt]{beamer}

% Shared style, ported from the UF Animal Sciences 2026 talk (ib-tutorials.sty at the repo root).
\makeatletter\def\input@path{{./}{../}}\makeatother
\usepackage{ib-tutorials}

\usepackage[utf8]{inputenc}
\usepackage{graphicx}
\usepackage{amsmath,amssymb,bm}
\usepackage{booktabs}
\usepackage{xcolor}
\usepackage{tcolorbox}
\tcbuselibrary{skins}
\usepackage{hyperref}

\graphicspath{{../pictures/}{../pictures/1-2/}{pictures/}{pictures/1-2/}}

\newcommand{\uvec}{\bm{u}}
\newcommand{\xvec}{\bm{x}}
\newcommand{\Xvec}{\bm{X}}
\newcommand{\Fvec}{\bm{F}}
\newcommand{\fvec}{\bm{f}}

\title[IB discretization]{Immersed boundary method discretization}
\subtitle{A power-consistent Eulerian--Lagrangian time step}
\author[Miller]{Laura A. Miller}
\institute[UArizona]{Department of Mathematics, University of Arizona}
\date{Fluid Dynamics and IBM Theory \textbullet\ Module 1}

\begin{document}

% ----- 1. Title -----
% Title slide via the shared \maketutorialtitlepanel macro so that the image
% can carry its source credit underneath (same layout as \maketutorialtitle).
\begin{frame}[plain]
\maketutorialtitlepanel{%
  \includegraphics[width=\linewidth,height=0.62\textheight,keepaspectratio]{ImmersedBoundary_grid2.png}\par\vspace{2pt}
  {\scriptsize Schematic after the Wikimedia Commons figure \emph{ImmersedBoundary grid} (Immersed boundary method article).\par}}
\end{frame}

% ----- 2. Goals -----
\begin{frame}{What this tutorial discretizes}
\begin{columns}[T,onlytextwidth]
\begin{column}{0.54\textwidth}
Starting from a continuum IB model, we need discrete operators for:
\begin{enumerate}\itemsep5pt
  \item incompressible flow on a fixed Cartesian grid;
  \item elastic forces on a moving material mesh;
  \item force spreading and velocity interpolation;
  \item advancement of the coupled configuration.
\end{enumerate}
\end{column}
\begin{column}{0.42\textwidth}
\begin{punchbox}
The formulas below define one \hi{first-order IMEX prototype}. They illustrate the coupling, not a unique IB time integrator.
\end{punchbox}
\medskip
\soft{IMEX: some terms are explicit and others implicit.}
\end{column}
\end{columns}
\end{frame}

% ----- 3. Codimension -----
\begin{frame}{One coupling idea spans several codimensions}
\begin{columns}[T,onlytextwidth]
\begin{column}{0.50\textwidth}
The immersed structure may be:
\begin{itemize}\itemsep5pt
  \item a one-dimensional fiber in a two- or three-dimensional fluid;
  \item a two-dimensional surface in a three-dimensional fluid;
  \item a volumetric solid represented by finite elements.
\end{itemize}
The interaction formulas keep the same form; only the material measure and its quadrature change.
\end{column}
\begin{column}{0.46\textwidth}
\begin{punchbox}
If $\mathcal U$ is the material domain,
\[
  \fvec(\xvec,t)=\int_{\mathcal U}\Fvec(q,t)
  \delta\bigl(\xvec-\boldsymbol\chi(q,t)\bigr)\,dq.
\]
Fibers, surfaces, and solids differ in the dimension and interpretation of $dq$.
\end{punchbox}
\end{column}
\end{columns}
\medskip
\soft{\small Notation: in this deck $\boldsymbol\chi(q,t)$ is the continuum current-position map, which decks 1-2 and 1-3 write as $\Xvec(q,t)$ (Peskin convention, material label $q$). Its discrete samples are $\Xvec_m^n=\boldsymbol\chi(q_m,t^n)$.}
\end{frame}

% ----- 4. Eulerian grid -----
\begin{frame}{Eulerian variables on a staggered grid}
For clarity, use a uniform two-dimensional MAC grid with spacing $h$:
\[
  p_{i,j}\approx p(ih,jh),\quad
  u_{i+\frac12,j}\approx u\!\left((i+\tfrac12)h,jh\right),\quad
  v_{i,j+\frac12}\approx v\!\left(ih,(j+\tfrac12)h\right).
\]
\begin{columns}[T,onlytextwidth]
\begin{column}{0.49\textwidth}
\begin{itemize}\itemsep4pt
  \item pressure is cell centered;
  \item each velocity component is face centered;
  \item the matching force component is stored on the same faces.
\end{itemize}
\end{column}
\begin{column}{0.47\textwidth}
\begin{punchbox}
The discrete gradient $G_h$ maps cells to faces, and the divergence $D_h$ maps faces to cells. With compatible boundary treatment, $D_h=-G_h^*$.
\end{punchbox}
\end{column}
\end{columns}
\end{frame}

% ----- 4. Lagrangian mesh -----
\begin{frame}{Lagrangian geometry and quadrature}
Let $q_m$ label material points and let
\[
  \Xvec_m^n=\boldsymbol{\chi}(q_m,t^n)
\]
denote their current positions. Associate each sample or quadrature point with a material weight $w_m$.
\begin{center}
\renewcommand{\arraystretch}{1.18}
\begin{tabular}{lll}
\toprule
Structure & Typical material measure & Weight \\
\midrule
fiber & $dq$ & $w_m\approx\Delta q$ \\
surface & $dA$ & surface quadrature weight \\
solid & $dV$ & volume quadrature weight \\
\bottomrule
\end{tabular}
\end{center}
\begin{punchbox}
$w_m$ is a \hi{quadrature measure}, not always a nodal spacing. In IBFE, quadrature points and structural nodes need not coincide.
\end{punchbox}
\end{frame}

% ----- 5. Interaction operators -----
\begin{frame}{Discrete interaction operators}
Using Eulerian sites $\xvec_{\bm{i}}$, cell/face weights $v_{\bm{i}}$, and a regularized kernel $\delta_h$, define
\begin{align*}
  (S_h[\Xvec]\Fvec)_{\bm{i}}
    &=\sum_m \Fvec_m\,\delta_h(\xvec_{\bm{i}}-\Xvec_m)\,w_m,
    &&\text{force spreading},\\
  (J_h[\Xvec]\uvec)_m
    &=\sum_{\bm{i}}\uvec_{\bm{i}}\,\delta_h(\xvec_{\bm{i}}-\Xvec_m)\,v_{\bm{i}},
    &&\text{velocity interpolation}.
\end{align*}
On a uniform Cartesian grid, $v_{\bm{i}}=h^d$; for a MAC grid, each component uses the locations of its own faces.
\begin{punchbox}
$S_h$ transfers a density per material measure to the fluid; $J_h$ samples the fluid velocity at material points.
\end{punchbox}
\end{frame}

% ----- 6. Adjointness -----
\begin{frame}{Matched coupling preserves discrete power}
Equip the two grids with their weighted inner products:
\[
  \langle \uvec,\fvec\rangle_E
  =\sum_{\bm{i}}\uvec_{\bm{i}}\!\cdot\!\fvec_{\bm{i}}v_{\bm{i}},
  \qquad
  \langle \bm{U},\Fvec\rangle_L
  =\sum_m \bm{U}_m\!\cdot\!\Fvec_mw_m.
\]
Substituting the definitions gives
\[
  \boxed{\langle \uvec,S_h\Fvec\rangle_E
  =\langle J_h\uvec,\Fvec\rangle_L},
  \qquad J_h=S_h^*.
\]
\begin{columns}[T,onlytextwidth]
\begin{column}{0.50\textwidth}
This is the discrete analogue of equal power exchanged between the fluid and structure.
\end{column}
\begin{column}{0.46\textwidth}
\begin{punchbox}
A common kernel is not enough: \hi{grids, weights, and inner products must be compatible}.
\end{punchbox}
\end{column}
\end{columns}
\end{frame}

% ----- 7. Classical implementation -----
\begin{frame}{From the general framework to a concrete scheme}
\begin{columns}[T,onlytextwidth]
\begin{column}{0.53\textwidth}
The coupling operators describe a \emph{family} of immersed boundary methods. To make the algorithm concrete, it is useful to examine the classical fiber-based construction associated with Peskin and McQueen's cardiac-flow work of the late 1980s and 1990s.
\medskip
\begin{itemize}\itemsep5pt
  \item The structure is a collection of material fibers.
  \item Elastic forces are assembled from targets, stretching, and bending.
  \item A regularized delta kernel couples the two grids.
\end{itemize}
\end{column}
\begin{column}{0.43\textwidth}
\begin{punchbox}
The following choices explain the logic of one influential implementation. They are \hi{not defining requirements} of every modern IB method.
\end{punchbox}
\end{column}
\end{columns}
\vfill
\soft{\scriptsize Adapted from: C.~S. Peskin and D.~M. McQueen, \emph{J. Comput. Phys.} 81 (1989) 372--405 (fiber-based heart model); C.~S. Peskin, \emph{Acta Numerica} 11 (2002) 479--517 (fiber forces, discretization, kernels).}
\end{frame}

\begin{frame}{Scope of the classical fiber discretization}
\small
The prototype considered here assumes:
\begin{itemize}\itemsep5pt
  \item an incompressible Newtonian fluid;
  \item a uniform Cartesian Eulerian grid;
  \item periodic boundaries when an FFT-based fluid solver is used;
  \item immersed fibers described in Lagrangian coordinates;
  \item viscosity treated implicitly, with advection and fluid--structure coupling treated explicitly.
\end{itemize}
\begin{punchbox}
Later IB and IBFE schemes relax these choices through physical boundary conditions, adaptive grids, continuum solid models, and semi-implicit or fully implicit coupling.
\end{punchbox}
\end{frame}

% ----- 9. Forces from energy -----
\begin{frame}{Compute elastic forces from a discrete energy}
If $E_h(\Xvec)$ is a discrete elastic energy, define force density by virtual work:
\[
  \delta E_h
  =-\sum_m \Fvec_m\cdot\delta\Xvec_m\,w_m,
  \qquad
  \Fvec_m=-\frac{1}{w_m}\frac{\partial E_h}{\partial\Xvec_m}.
\]
\begin{columns}[T,onlytextwidth]
\begin{column}{0.50\textwidth}
Common contributions include:
\begin{itemize}\itemsep4pt
  \item stretching or area/volume elasticity;
  \item bending and preferred curvature;
  \item active stresses or muscle models;
  \item target-point penalties.
\end{itemize}
\end{column}
\begin{column}{0.46\textwidth}
For a prescribed target $\bm{Y}_m^n$,
\[
  \Fvec_{\mathrm{targ},m}^n
  =k_{\mathrm{targ}}(\bm{Y}_m^n-\Xvec_m^n).
\]
\soft{Large penalty stiffness usually makes an explicit scheme more restrictive.}
\end{column}
\end{columns}
\end{frame}

% ----- 8. Stretching -----
\begin{frame}{Example: conservative fiber tension}
For equally spaced material labels, put tension on links:
\[
  \bm{d}_{m+\frac12}=\frac{\Xvec_{m+1}-\Xvec_m}{\Delta q},
  \qquad
  \bm{\tau}_{m+\frac12}=\frac{\bm{d}_{m+\frac12}}{|\bm{d}_{m+\frac12}|},
\]
\[
  \bm{t}_{m+\frac12}=T\!\left(|\bm{d}_{m+\frac12}|\right)\bm{\tau}_{m+\frac12},
  \qquad
  \boxed{\Fvec_m^{\mathrm{str}}
  =\frac{\bm{t}_{m+\frac12}-\bm{t}_{m-\frac12}}{\Delta q}}.
\]
\begin{punchbox}
The flux-difference form makes internal link forces equal and opposite. The constitutive law $T(\lambda)$ sets the reference stretch and stiffness.
\end{punchbox}
\vspace{3pt}
\soft{Endpoint forces require the physical boundary conditions for the fiber.}
\end{frame}

% ----- 9. Bending -----
\begin{frame}{Example: bending is a restoring force}
For the energy
\[
  E_b=\frac{\kappa_b}{2}\int
  |\boldsymbol{\chi}_{qq}-\boldsymbol{\chi}^{\,0}_{qq}|^2\,dq,
\]
the interior force is
\[
  \Fvec_b=-\kappa_b\,\partial_q^4
  (\boldsymbol{\chi}-\boldsymbol{\chi}^{\,0}),
\]
up to endpoint terms fixed by the boundary conditions. For a straight preferred configuration,
\[
  \Fvec_{b,m}\approx-\kappa_b
  \frac{\Xvec_{m-2}-4\Xvec_{m-1}+6\Xvec_m-4\Xvec_{m+1}+\Xvec_{m+2}}
       {(\Delta q)^4}.
\]
\begin{punchbox}
The minus sign is essential: the elastic force decreases the bending energy.
\end{punchbox}
\end{frame}

\begin{frame}{Classical fiber forces: three additive pieces}
\small
For material nodes $\Xvec_m^n\approx\boldsymbol\chi(q_m,t^n)$ with spacing $\Delta q$, a common decomposition is
\[
 \Fvec_m^n=\Fvec_{\rm targ,m}^n+\Fvec_{\rm str,m}^n+\Fvec_{\rm bend,m}^n.
\]
Target points give the penalty force
\[
 \Fvec_{\rm targ,m}^n=k_{\rm targ}(\bm Y_m^n-\Xvec_m^n).
\]
With
\[
 (D_q\Xvec)_m=\frac{\Xvec_{m+1}-\Xvec_{m-1}}{2\Delta q},
 \qquad \bm\tau_m=\frac{(D_q\Xvec)_m}{|(D_q\Xvec)_m|},
\]
one illustrative Hookean stretching law is
\[
 \Fvec_{\rm str,m}=k_{\rm str}D_q\!\left[(|D_q\Xvec|-1)\bm\tau\right]_m.
\]
\soft{The link-tension form shown earlier is often preferable for making pairwise force balance transparent.}
\end{frame}

\begin{frame}{Classical fiber bending: sign and boundaries}
For a straight preferred configuration, the interior bending force is
\[
 \Fvec_{\rm bend,m}=-k_{\rm bend}
 \frac{\Xvec_{m-2}-4\Xvec_{m-1}+6\Xvec_m-4\Xvec_{m+1}+\Xvec_{m+2}}
      {(\Delta q)^4}.
\]
\begin{columns}[T,onlytextwidth]
\begin{column}{0.50\textwidth}
The five-point stencil approximates $-k_{\rm bend}\partial_q^4\boldsymbol\chi$ at interior nodes.
\end{column}
\begin{column}{0.46\textwidth}
\begin{punchbox}
Closed fibers use periodic indexing. Open fibers require endpoint conditions consistent with the underlying bending energy.
\end{punchbox}
\end{column}
\end{columns}
\medskip
\warn{Sign check:} positive $k_{\rm bend}$ must produce a restoring force, so the discrete operator carries the leading minus sign.
\end{frame}

% ----- 10. Fluid solve -----
\begin{frame}{A first-order IMEX fluid step}
One pedagogical prototype treats advection and coupling explicitly and viscosity implicitly:
\begin{align*}
  \rho\left(\frac{\uvec^{n+1}-\uvec^n}{\Delta t}
    +N_h(\uvec^n)\right)
  &=-G_hp^{n+1}+\mu L_h\uvec^{n+1}
    +S_h[\Xvec^n]\Fvec^n+\bm{g}^n,\\
  D_h\uvec^{n+1}&=0.
\end{align*}
Here $L_h$ is the face-centered vector Laplacian and $N_h$ is the chosen discrete convection operator.
\begin{punchbox}
Pressure appears as $-G_hp$ on the right-hand side (equivalently, $+G_hp$ on the left). The opposite sign would reverse the pressure force.
\end{punchbox}
\end{frame}

\begin{frame}{A component form of the classical fluid update}
\small
On a uniform grid, the same prototype may be written with finite-difference operators as
\begin{align*}
 \rho\left[
   \frac{\uvec^{n+1}-\uvec^n}{\Delta t}
   +\mathcal S_h(\uvec^n)\uvec^n
 \right]
 &=-\bm D^{0}p^{n+1}
   +\mu\sum_{\alpha=1}^{d}D_\alpha^+D_\alpha^-\uvec^{n+1}
   +\fvec^n,\\
 \bm D^{0}\!\cdot\uvec^{n+1}&=0.
\end{align*}
The term $-\bm D^{0}p^{n+1}$ here and the term $-G_hp^{n+1}$ on the previous slide are \emph{two discretizations of the same} continuum pressure-gradient term $-\nabla p$, written on two grids: $G_h$ is the staggered (MAC) gradient, $\bm D^{0}$ its collocated centered-difference counterpart (the form used in Peskin 2002). They are not the same matrix: the collocated $\bm D^{0}$ gradient has a $2^d$-dimensional null space (constants on each even/odd ``chain'' of nodes), which is why Peskin's four-point kernel imposes the even--odd condition (deck 1-2).
Once $\fvec^n$ and the explicit convection term are known, the unknown velocity and pressure enter linearly.
\begin{punchbox}
On a periodic box, this structure makes FFT-based solvers natural. With physical boundaries or local refinement, other linear solvers are used.
\end{punchbox}
\end{frame}

\begin{frame}{Finite-difference operators in that update}
\small
For a scalar grid function $\phi$ and Cartesian unit vector $\bm e_\alpha$,
\begin{align*}
 D_\alpha^0\phi(\xvec)
 &=\frac{\phi(\xvec+h\bm e_\alpha)-\phi(\xvec-h\bm e_\alpha)}{2h},\\
 D_\alpha^+\phi(\xvec)
 &=\frac{\phi(\xvec+h\bm e_\alpha)-\phi(\xvec)}{h},\qquad
 D_\alpha^-\phi(\xvec)
 =\frac{\phi(\xvec)-\phi(\xvec-h\bm e_\alpha)}{h}.
\end{align*}
One energy-conscious convection form is
\[
 \mathcal S_h(\uvec)\phi
 =\tfrac12\uvec\cdot\bm D^0\phi
 +\tfrac12\bm D^0\!\cdot(\uvec\phi).
\]
\begin{punchbox}
On a staggered grid, the stencils differ; the design goal is unchanged: a compatible gradient/divergence pair with controlled kinetic-energy behavior.
\end{punchbox}
\end{frame}

% ----- Projection view of the fluid step (added: addresses the open audit item) -----
\begin{frame}{Solving the fluid step: a projection (fractional-step) view}
\small
\setlength{\abovedisplayskip}{4pt}\setlength{\belowdisplayskip}{4pt}%
A common way to solve the IMEX system splits it into a viscous step and a projection onto discretely divergence-free fields:
\begin{align*}
  &\text{1. viscous predictor:} &
  \rho\,\frac{\tilde{\uvec}-\uvec^n}{\Delta t}+\rho N_h(\uvec^n)
  &=\mu L_h\tilde{\uvec}+S_h[\Xvec^n]\Fvec^n,\\
  &\text{2. pressure Poisson:} &
  D_hG_h\,\varphi&=\frac{\rho}{\Delta t}\,D_h\tilde{\uvec},\\
  &\text{3. projection:} &
  \uvec^{n+1}&=\tilde{\uvec}-\frac{\Delta t}{\rho}\,G_h\varphi
  \quad\Rightarrow\quad D_h\uvec^{n+1}=0 .
\end{align*}
\begin{columns}[T,onlytextwidth]
\begin{column}{0.50\textwidth}
\textbf{Periodic, uniform grid.} \(L_h\), \(G_h\), \(D_h\) commute, and substituting step 3 into step 1 reproduces the coupled IMEX system exactly with \(p^{n+1}=(I-\nu\Delta t\,L_h)\varphi\) (\(L_h\) acting as the scalar Laplacian), \(\nu=\mu/\rho\). FFTs diagonalize every step.
\end{column}
\begin{column}{0.46\textwidth}
\textbf{Physical boundaries.} The operators no longer commute; the split introduces a splitting error (often limiting pressure accuracy near boundaries) and needs boundary conditions for \(\tilde{\uvec}\) and \(\varphi\). Coupled (Stokes-type) solvers avoid the split.
\end{column}
\end{columns}
\vspace{2pt}
\soft{\scriptsize After A.~J. Chorin, \emph{Math. Comp.} 22 (1968) 745--762; D.~L. Brown, R. Cortez, M.~L. Minion, \emph{J. Comput. Phys.} 168 (2001) 464--499.}
\end{frame}

% ----- 11. One time step -----
\begin{frame}{One consistent explicit-coupling time step}
Given $(\uvec^n,\Xvec^n)$:
\begin{enumerate}\itemsep6pt
  \item compute $\Fvec^n=\mathcal{F}_h(\Xvec^n,t^n)$;
  \item spread $\fvec^n=S_h[\Xvec^n]\Fvec^n$;
  \item solve the incompressible fluid step for $(\uvec^{n+1},p^{n+1})$;
  \item interpolate $\bm{U}^{n+1}=J_h[\Xvec^n]\uvec^{n+1}$;
  \item update $\Xvec^{n+1}=\Xvec^n+\Delta t\,\bm{U}^{n+1}$.
\end{enumerate}
\begin{punchbox}
Every occurrence of the configuration in this prototype is evaluated at $\Xvec^n$. Higher-order, midpoint, split, and implicit schemes use different---but deliberately consistent---time levels.
\end{punchbox}
\end{frame}

\begin{frame}{No slip is enforced through interpolation}
At a material point $q_m$, the continuum kinematic equation is regularized as
\[
 \frac{d\boldsymbol\chi}{dt}(q_m,t)
 \approx\int_\Omega \uvec(\xvec,t)\,
 \delta_h\bigl(\xvec-\boldsymbol\chi(q_m,t)\bigr)\,d\xvec.
\]
Its discrete counterpart is $\bm U_m=J_h[\Xvec]\uvec$.
\begin{columns}[T,onlytextwidth]
\begin{column}{0.50\textwidth}
\begin{itemize}\itemsep4pt
  \item the fluid velocity is averaged over the kernel support;
  \item the structure advances with that interpolated velocity;
  \item the same matched kernel spreads structural force back to the fluid.
\end{itemize}
\end{column}
\begin{column}{0.46\textwidth}
\begin{punchbox}
At finite $h$, this is a regularized or weak no-slip relation rather than pointwise sampling of an exact interface trace.
\end{punchbox}
\end{column}
\end{columns}
\end{frame}

% ----- 12. Stiffness -----
\begin{frame}{Why bending and penalties can be stiff}
For the discrete fourth derivative,
\[
  \lambda_{\max}(D_q^4)=\mathcal{O}((\Delta q)^{-4}).
\]
Thus mesh-scale bending modes become rapidly stiffer as the Lagrangian mesh is refined.
\begin{columns}[T,onlytextwidth]
\begin{column}{0.48\textwidth}
\hi{What this does mean}
\begin{itemize}\itemsep4pt
  \item explicit elastic coupling may impose a severe $\Delta t$ restriction;
  \item large target stiffness can have the same effect.
\end{itemize}
\end{column}
\begin{column}{0.48\textwidth}
\hi{What it does not mean}
\begin{itemize}\itemsep4pt
  \item a resolved physical bending force itself grows like $(\Delta q)^{-4}$;
  \item IBFE or a penalty law automatically removes stiffness.
\end{itemize}
\end{column}
\end{columns}
\end{frame}

% ----- 13. Stability -----
\begin{frame}{Time-step restrictions are scheme dependent}
\small
\renewcommand{\arraystretch}{1.22}
\begin{tabular}{@{}p{0.22\textwidth}p{0.34\textwidth}p{0.32\textwidth}@{}}
\toprule
Explicit term & Representative scale & Important dependence \\
\midrule
advection & $\Delta t\,\max\limits_{\bm{i}}\sum_\alpha |u_\alpha|/h_\alpha\lesssim C_{\rm adv}$ & flux and time integrator \\
viscous diffusion & $\nu\Delta t/h^2\lesssim C_\nu$ & none with implicit diffusion \\
elastic coupling & no universal one-line bound & force law, grid, kernel, solver \\
\bottomrule
\end{tabular}
\medskip
\begin{punchbox}
There is no general rule that $\Delta t$ is simply ``inversely proportional to the applied force.'' Diagnose the unstable mode and the terms treated explicitly.
\end{punchbox}
\end{frame}

% ----- 14. Resolution -----
\begin{frame}{Resolve the interaction, not just the geometry}
\begin{columns}[T,onlytextwidth]
\begin{column}{0.51\textwidth}
\begin{itemize}\itemsep5pt
  \item \warn{Too coarse:} gaps between interaction stencils can cause leakage and volume error.
  \item \warn{Unnecessarily fine:} more work, redundant constraints, and possibly worse conditioning or high-frequency stiffness.
  \item \hi{IBFE:} structural interpolation and interaction quadrature can be refined independently.
\end{itemize}
\end{column}
\begin{column}{0.45\textwidth}
\begin{punchbox}
Rules such as $\Delta q\approx h/2$ are method- and kernel-dependent heuristics, not universal laws. Check leakage and convergence for the chosen discretization.
\end{punchbox}
\medskip
The kernel support, geometry, quadrature rule, and local Eulerian spacing all matter.
\end{column}
\end{columns}
\end{frame}

% ----- 15. Kernel -----
\begin{frame}{The kernel sets the interaction stencil}
A common separable construction is
\[
  \delta_h(\xvec)=\frac{1}{h^d}\prod_{\alpha=1}^{d}
  \phi\!\left(\frac{x_\alpha}{h}\right).
\]
Useful discrete identities include
\[
  \sum_{j\in\mathbb{Z}}\phi(r-j)=1,
  \qquad
  \sum_{j\in\mathbb{Z}}(r-j)\phi(r-j)=0.
\]
\begin{columns}[T,onlytextwidth]
\begin{column}{0.50\textwidth}
The zeroth moment preserves total force and interpolates constants; the first moment centers spreading and supports torque/linear-field reproduction.
\end{column}
\begin{column}{0.46\textwidth}
\begin{punchbox}
Tensor-product kernels are efficient and grid aligned; they are not exactly invariant under rotation.
\end{punchbox}
\end{column}
\end{columns}
\end{frame}

\begin{frame}{Example: a four-point-support cosine kernel}
One simple compactly supported choice is
\[
 \phi(r)=
 \begin{cases}
  \dfrac14\left(1+\cos\dfrac{\pi r}{2}\right), & |r|\le 2,\\[6pt]
  0, & |r|>2,
 \end{cases}
 \qquad
 \delta_h(\xvec)=\frac{1}{h^d}\prod_{\alpha=1}^{d}\phi(x_\alpha/h).
\]
\begin{columns}[T,onlytextwidth]
\begin{column}{0.50\textwidth}
It is symmetric and compactly supported: in each coordinate, only grid sites within $2h$ contribute.\par\smallskip
\soft{\scriptsize It is Peskin's simple approximation to his four-point kernel (Peskin 2002, eq.~6.28), and it does satisfy the even--odd condition.}
\end{column}
\begin{column}{0.46\textwidth}
\begin{punchbox}
``Four-point support'' describes the one-dimensional stencil width. In $d$ dimensions, the tensor-product support contains up to $4^d$ grid sites per Lagrangian point.
\end{punchbox}
\end{column}
\end{columns}
\smallskip
\warn{Note:} $\phi_{\cos}$ satisfies the zeroth moment and $\sum_j\phi_{\cos}(r-j)^2=3/8$ but \emph{not} the first-moment condition (max $|M_1|\approx0.021$); see deck 1-3.\par
\soft{\scriptsize Here $M_1(r)=\sum_j(r-j)\phi_{\cos}(r-j)=r-\tfrac12+\tfrac12\bigl(\cos\tfrac{\pi r}{2}-\sin\tfrac{\pi r}{2}\bigr)$ for $0\le r<1$, in units of $h$; the maximum occurs at $r\approx0.21$ and $0.79$.}
\end{frame}

% ----- 16. Verification -----
\begin{frame}{Verification checks for an implementation}
\begin{columns}[T,onlytextwidth]
\begin{column}{0.49\textwidth}
\hi{Operator checks}
\begin{itemize}\itemsep4pt
  \item $D_h=-G_h^*$ under the chosen boundary treatment;
  \item $J_h=S_h^*$ in weighted inner products;
  \item constant and linear-field kernel tests;
  \item total force and torque after spreading.
\end{itemize}
\end{column}
\begin{column}{0.47\textwidth}
\hi{Solution checks}
\begin{itemize}\itemsep4pt
  \item incompressibility and solver residuals;
  \item energy or power balance when applicable;
  \item independent refinement in $h$, quadrature/mesh spacing, and $\Delta t$;
  \item sensitivity to domain size and boundary conditions.
\end{itemize}
\end{column}
\end{columns}
\begin{punchbox}
Refining every discretization parameter together can hide the source of an error.
\end{punchbox}
\end{frame}

% ----- 17. Accuracy -----
\begin{frame}{What convergence requires}
For a conventional regularized-delta IB discretization:
\begin{itemize}\itemsep5pt
  \item the interaction layer has width $\mathcal{O}(h)$;
  \item singular forces and pressure/traction jumps are therefore smoothed at finite $h$;
  \item interfacial quantities are often less accurate than smooth fields away from the interface.
\end{itemize}
\begin{punchbox}
$\delta_h\to\delta$ distributionally is \warn{not by itself} a convergence proof. The fluid scheme, time integrator, material mesh or quadrature, and coupling must be stable and refined consistently.
\end{punchbox}
\end{frame}

% ----- 18. Extensions -----
\begin{frame}{From the prototype to modern IB solvers}
The same interaction-operator viewpoint extends to:
\begin{itemize}\itemsep5pt
  \item block-structured adaptive mesh refinement;
  \item semi-implicit or fully implicit fluid--structure coupling;
  \item finite-element descriptions of volumetric solids (IBFE);
  \item constrained-body and sharp-interface variants;
  \item scalable multilevel linear and nonlinear solvers.
\end{itemize}
\begin{punchbox}
Uniform grids, periodic boxes, FFT solvers, and explicit coupling are choices of particular implementations---not defining assumptions of the immersed boundary formulation.
\end{punchbox}
\end{frame}

\begin{frame}{Classical choices and modern extensions}
\begin{columns}[T,onlytextwidth]
\begin{column}{0.46\textwidth}
\begin{punchbox}
\centering\hi{Classical fiber prototype}
\begin{itemize}\itemsep4pt
  \item fiber structures;
  \item spring and beam elasticity;
  \item commonly explicit coupling;
  \item uniform, often periodic grid;
  \item FFT solver on a periodic box.
\end{itemize}
\end{punchbox}
\end{column}
\begin{column}{0.50\textwidth}
\begin{punchbox}
\centering\hi{Modern IB/IBFE examples}
\begin{itemize}\itemsep4pt
  \item fibers, surfaces, and volumetric solids;
  \item FE continuum and active-material laws;
  \item explicit through monolithic coupling;
  \item physical boundaries and adaptive grids;
  \item scalable multigrid and Krylov solvers.
\end{itemize}
\end{punchbox}
\end{column}
\end{columns}
\medskip
\begin{punchbox}
The enduring idea is the Eulerian--Lagrangian interaction pair. The structure model, grid hierarchy, and time integrator can all change around it.
\end{punchbox}
\end{frame}

\begin{frame}{Roadmap from the classical method}
\begin{columns}[T,onlytextwidth]
\begin{column}{0.47\textwidth}
\hi{What the prototype teaches}
\begin{itemize}\itemsep4pt
  \item moving geometry without remeshing;
  \item forces and kinematics transferred by a common kernel;
  \item a transparent sequence for one coupled time step.
\end{itemize}
\end{column}
\begin{column}{0.49\textwidth}
\hi{What motivates later methods}
\begin{itemize}\itemsep4pt
  \item stiffness under explicit elastic coupling;
  \item reduced interfacial accuracy from regularization;
  \item the need for continuum solids, adaptivity, and robust solvers.
\end{itemize}
\end{column}
\end{columns}
\begin{punchbox}
These limitations lead naturally to implicit IB schemes, immersed interface variants, adaptive solvers, and the IBFE formulation developed later in this tutorial series.
\end{punchbox}
\end{frame}

% ----- 19. Summary -----
\begin{frame}{Discretization checklist}
\begin{enumerate}\itemsep7pt
  \item Put Eulerian variables and discrete differential operators on compatible locations.
  \item Treat material weights as quadrature measures and compute forces with the correct sign.
  \item Construct $S_h$ and $J_h$ as weighted adjoints so coupling preserves power.
  \item State the time levels and implicit/explicit split unambiguously.
  \item Test conservation, leakage, stability, and convergence independently.
\end{enumerate}
\end{frame}

% ----- 20. References -----
\begin{frame}{References and further study}
\begin{itemize}\itemsep6pt
  \item C. S. Peskin, ``The immersed boundary method,'' \emph{Acta Numerica} 11 (2002),
  \href{https://doi.org/10.1017/S0962492902000077}{doi:10.1017/S0962492902000077}.
  \item B. E. Griffith and X. Y. Luo, ``Hybrid finite difference/finite element immersed boundary method,'' \emph{Int. J. Numer. Meth. Biomed. Engng.} 33 (2017), e2888,
  \href{https://doi.org/10.1002/cnm.2888}{doi:10.1002/cnm.2888}.
  \item B. E. Griffith and N. A. Patankar, ``Immersed methods for fluid--structure interaction,'' \emph{Annu. Rev. Fluid Mech.} 52 (2020),
  \href{https://doi.org/10.1146/annurev-fluid-010719-060228}{doi:10.1146/annurev-fluid-010719-060228}.
  \item A. M. Roma, C. S. Peskin, and M. J. Berger, ``An adaptive version of the immersed boundary method,'' \emph{J. Comput. Phys.} 153 (1999), 509--534.
  \item \href{https://ibamr.github.io/}{IBAMR documentation and examples}.
\end{itemize}
\end{frame}

\begin{frame}{References and further study (continued)}
\footnotesize
\begin{itemize}\itemsep3pt
  \item C.~S. Peskin, ``Numerical analysis of blood flow in the heart,'' \emph{J. Comput. Phys.} 25 (1977), 220--252,
  \href{https://doi.org/10.1016/0021-9991(77)90100-0}{doi:10.1016/0021-9991(77)90100-0}.
  \item C.~S. Peskin and D.~M. McQueen, ``A three-dimensional computational method for blood flow in the heart. I. Immersed elastic fibers in a viscous incompressible fluid,'' \emph{J. Comput. Phys.} 81 (1989), 372--405.
  \item F.~H. Harlow and J.~E. Welch, ``Numerical calculation of time-dependent viscous incompressible flow of fluid with free surface,'' \emph{Phys. Fluids} 8 (1965), 2182--2189 (MAC grid),
  \href{https://doi.org/10.1063/1.1761178}{doi:10.1063/1.1761178}.
  \item B.~E. Griffith, R.~D. Hornung, D.~M. McQueen, and C.~S. Peskin, ``An adaptive, formally second order accurate version of the immersed boundary method,'' \emph{J. Comput. Phys.} 223 (2007), 10--49.
  \item A.~J. Chorin, ``Numerical solution of the Navier--Stokes equations,'' \emph{Math. Comp.} 22 (1968), 745--762 (projection method).
  \item D.~L. Brown, R. Cortez, and M.~L. Minion, ``Accurate projection methods for the incompressible Navier--Stokes equations,'' \emph{J. Comput. Phys.} 168 (2001), 464--499.
  \item Title figure: schematic after the Wikimedia Commons figure \emph{ImmersedBoundary grid} (Immersed boundary method article).
\end{itemize}
\end{frame}

\end{document}
